% Unregistered candidate D.59.  Root audit requested before integration.
\paragraph{Candidate D.59: every second prime detects the Fermat family.}

Let $\ell\in\{3,5\}$, $n=\ell^r$, $k=\mathbb Q_\ell$,
and use the Fermat curve, cusp and family of D.20/D.57:
\[
 X_n:X^n+Y^n+Z^n=0,\quad
 U_n=X_n\setminus\{XYZ=0\},\quad O=[0:-1:1],\quad
 s_a\quad(a\in\mathbb Z_\ell^{d_n}).
\]
Write $J_n$ for its Jacobian.  The single-prime proper Kummer map
\[
 X_n(k)\longrightarrow H^1(k,T_\ell J_n),\qquad
                   P\longmapsto\kappa_\ell([P-O])       \tag{D.59.1}
\]
is injective.  In particular its zero fibre is $\{O\}$.

For every prime $q\ne\ell$, put $\Sigma=\{\ell,q\}$ and
$N_\Sigma=\pi_1(U_{n,\overline{\mathbb Q}})^\Sigma$.
The images of the $s_a$ in $N_\Sigma/N_\Sigma''$ remain
pairwise nonconjugate over $k$.  Every nonzero parameter
has no local lift to $N_\Sigma$, and hence no global lift.
For each $a\ne0$ there is a finite $m(a,q)\ge3$ such that
it already has no local lift to
$N_\Sigma/N_\Sigma^{(m(a,q))}$.  For each $h\ge1$ the
stage can be chosen uniformly for
$a\notin\ell^h\mathbb Z_\ell^{d_n}$.
Thus $q=2$ suffices in both the three-power and five-power
towers.  The least obstructing derived length and an
explicit finite obstruction cocycle remain undetermined.

\begin{proof}
\paragraph{1. All local points lie in three explicitly controlled disks.}

For a unit modulo $\ell^2$, its $n$th power equals its
$\ell$th power, because
\[
 n\equiv\ell\pmod{\ell(\ell-1)}.
\]
The possible unit powers are
\[
 \begin{array}{c|c|c}
 \ell&\ell^2&\{u^n\bmod\ell^2:u\in\mathbb Z_\ell^\times\}\\\hline
 3&9&\{1,8\}\\
 5&25&\{1,7,18,24\}.
 \end{array}                                          \tag{D.59.2}
\]
No three entries in either set, with repetition allowed,
sum to zero modulo $\ell^2$.  This finite assertion is
executed exhaustively in \path{code/de_fermat_two_prime.py}:
there are eight triples for $\ell=3$ and 64 for $\ell=5$.
The same code checks every primitive triple modulo
$\ell^2$ satisfying the Fermat equation; all 162 triples
modulo nine and all 1500 triples modulo 25 have exactly
one coordinate divisible by $\ell$.

Normalize a point of $X_n(k)$ to primitive integral
projective coordinates.  Equation (D.59.2) rules out
three unit coordinates.  Two coordinates divisible by
$\ell$ would leave the remaining unit term nonzero
modulo $\ell$, which is also impossible.  Therefore
exactly one coordinate is divisible by $\ell$.

Suppose that coordinate is $X$.  With $t=X/Z$ and
$u=Y/Z$ we have $t\in\ell\mathbb Z_\ell$, $u$ a unit,
and $u^n=-(1+t^n)$.  Reduction modulo $\ell$ gives
$u\equiv-1\pmod\ell$.  The analytic expression
\[
 h(t)=\exp\left(\frac1n\log(1+t^n)\right)              \tag{D.59.3}
\]
converges on the geometric open disk
\[
 v_\ell(t)>\rho_{\ell,r}:=
                    \frac{r+1/(\ell-1)}{\ell^r}<1.
\]
Indeed $v_\ell(t^n)>r+1/(\ell-1)$, so the logarithm
has that same valuation, and division by $n$ leaves
valuation greater than $1/(\ell-1)$, the convergence
threshold for the exponential.  The strict inequality
$\rho_{\ell,r}<1$ holds at $r=1$ and then follows by
induction.  Thus (D.59.3) is analytic on a disk strictly
larger than the closed disk $|t|\le\ell^{-1}$.

It satisfies $h(t)^n=1+t^n$ and $h(0)=1$.
The group $1+\ell\mathbb Z_\ell$ is torsion free for
$\ell=3,5$, as its logarithm injects it into
$\ell\mathbb Z_\ell$.  Hence $u=-h(t)$ is the unique
$k$-valued branch reducing to $-1$.  The map
\[
 t\longmapsto[t:-h(t):1]
\]
is an analytic embedding of that disk in $X_n^{\rm an}$:
projection to $t$ is its inverse, and the derivative
$nu^{n-1}$ in the defining equation is nonzero.  It
contains $O$ and every local point for which $X$ is
the nonunit coordinate.  Permuting coordinates gives
the other two disks, centred at
\[
                      O_1=[-1:0:1],\qquad
                      O_\infty=[-1:1:0].
\]
Their closed subdisks $|t|\le\ell^{-1}$ cover $X_n(k)$.
Only these embedded closed disks are needed below.

\paragraph{2. Prime-to-$\ell$ Kummer classes are constant on each disk.}

Fix a prime $b\ne\ell$, an integer $j\ge1$, and one
of the three disk centres $d$.  Let
$\iota_d:X_n\to J_n$ be $P\mapsto[P-d]$.  Pull back
$[b^j]:J_n\to J_n$ along $\iota_d$.  This is a finite
etale torsor under $J_n[b^j]$, including at the cusp $d$.
After base change to $\mathbb C_\ell$ it is a Galois
cover of degree a power of $b$.  Its restriction to
the embedded closed disk is trivial by Berkovich's
Theorem 6.3.2 and Corollary 6.3.3, printed pages
119--121 \cite{DE_Berkovich1993}.  The exact scope is
prime-to-residue-characteristic geometric etale
covers; no triviality assertion for $\ell$-power
covers of a disk is used.

The trivialization gives a canonical identification
of the fibres over $d$ and over any $P\in X_n(k)$
in the disk: pair points in the same connected
component of the covering disk.  This identification
commutes with the action of $G_k$, since the disk,
covering and its centre are defined over $k$.
The rational zero division point in the fibre over
$d$ belongs to a unique component.  That component
is $G_k$-stable, so its unique point over $P$ is
$G_k$-fixed.  This point belongs to the finite
algebraic fibre over $P$, hence is a $k$-rational
division point, not merely a point over a completed
extension.  It follows that
\[
                     [P-d]\in b^jJ_n(k)
                       \quad\text{for every }j.       \tag{D.59.4}
\]
This proves the vanishing of its full $b$-adic
Kummer class with the required rational descent.

The line $Y+Z=0$ cuts the Fermat curve in $nO$.
Similarly $X+Z=0$ and $X+Y=0$ cut it in $nO_1$
and $nO_\infty$.  Ratios of these lines therefore
show
\[
                   n[O_1-O]=n[O_\infty-O]=0.
\]
These differences are $\ell$-primary torsion and
are divisible by every power of $b\ne\ell$ inside
$J_n(k)$.  Combining this with (D.59.4) gives
\[
 \kappa_b([P-O])=0\quad
       \text{for every }P\in X_n(k),\ b\ne\ell.         \tag{D.59.5}
\]

If two points have the same class in (D.59.1),
their difference has zero Kummer class at every
prime, by (D.59.5).  The full local Kummer map of
$J_n(k)$ is injective by the profinite argument
in D.57.  Thus their degree-zero divisor classes
coincide.  The Abel--Jacobi map of $X_n$, whose
genus is at least one, is injective on points.
The two points are equal, proving (D.59.1).

\paragraph{3. Projection to any two-prime metabelian group retains the parameters.}

The geometric quotient map
$N_\Sigma\twoheadrightarrow N^{(\ell)}=P$
induces maps on abelianization, commutator
abelianization, and the class-two quotient
$B=\bigwedge^2P^{\rm ab}$ used in D.57.
Consequently the projection of the actual D.20
cocycle still has parameter $a$ in
$H^1(G_k,B)$.  This parameter is nonzero when
$a\ne0$, by the actual local surjection to
$\Gamma_\ell$ and inflation injectivity proved
in D.57.

Set $A_\Sigma=N_\Sigma^{\rm ab}$.  The proof
of (D.57.2), restricted to the two coefficient
primes, gives
\[
 H^0(G_k,A_\Sigma)=0,\qquad
 H^0(G_k,A_\Sigma/I_{O,\Sigma})=0.                    \tag{D.59.6}
\]
Here $I_{O,\Sigma}=\prod_{b\in\Sigma}\mathbb Z_b(1)$
is the primitive cusp line.  Thus a conjugator
between two projected family sections has zero
abelian image and lies in the commutator kernel.
Its projection to $B$ would make their difference
parameter a coboundary.  The two parameters must
therefore be equal.  This proves pairwise
nonconjugacy in the two-prime metabelian quotient.

Every family section has zero abelian cocycle.
If its local two-prime class were geometric, its
proper $\ell$-adic Kummer class would be zero.
By (D.59.1) the underlying compactification point
would be $O$.  Hence it would be cuspidal at
$O$.  The second vanishing in (D.59.6) makes
$H^1(k,I_{O,\Sigma})\to H^1(k,A_\Sigma)$
injective.  The zero abelian class forces the
unit cusp parameter, whose metabelian class is
$s_0$.  Pairwise nonconjugacy excludes this
for $a\ne0$.

\paragraph{4. A two-prime local lift has to be geometric.}

The class of all finite groups whose orders
have prime divisors in $\Sigma$ is a full
formation.  It contains the residue prime
$\ell$ and the distinct prime $q$.  D.56
therefore makes every local section of the
geometrically pro-$\Sigma$ arithmetic group
point or cusp induced.  Paragraph 3 excludes
a lift of every nonzero family member to
that group.  Restriction to $G_k$ excludes
a global lift as well.

Every finite $\Sigma$-group is solvable by
Burnside's two-prime theorem.  Therefore
\[
 N_\Sigma=\varprojlim_m
                 N_\Sigma/N_\Sigma^{(m)}.
\]
The actual local cocycle spaces for its
finite derived quotients are compact by
the finite-solvable-coefficient argument
of D.57.  If a fixed metabelian class
lifted at every derived length, conjugator
lifting makes all fibres exact, and
compactness gives one compatible lift to
$N_\Sigma$.  This is impossible for a
nonzero parameter.  The closed nested
liftable-parameter subsets and compactness
of the complement of
$\ell^h\mathbb Z_\ell^{d_n}$ give the
uniform finite stage exactly as in
(D.57.4).  No surjectivity of transition
maps of section fibres is assumed.
\end{proof}
